\begin{corollary}\label{sga3:XXI:4.3}
\footnote{Editors' note: The statement was corrected.}
Let \(x\in V^\ast\). For every \(w\in W\), put
\[
 x-w(x)=a_w\alpha^\ast+b_w\beta^\ast.
\]
If \((x,\alpha)\ge0\) and \((x,\beta)\ge0\), then
\[
 a_w+b_w\ge0.
\]
If moreover \((x,\alpha)>0\) (respectively \((x,\beta)>0\)), then
\[
 a_w+b_w>0
\]
for \(w\ne1,s_\beta\) (respectively for \(w\ne1,s_\alpha\)).
\end{corollary}

\begin{proof}
This follows at once from Proposition~\ref{sga3:XXI:4.1} and
Lemma~\ref{sga3:XXI:4.2}.
\end{proof}

\begin{corollary}\label{sga3:XXI:4.4}
Let \(\mathcal{R}\) be any root datum, and let \(\Delta\) be a system
of simple roots. Let \(\gamma\) be a positive root, let
\(\alpha,\beta\) be two simple roots, and let \(W_{\alpha,\beta}\) be
the subgroup of \(W\) generated by \(s_\alpha\) and \(s_\beta\). If
\[
 \operatorname{ord}_\Delta(s_\alpha(\gamma))
 <\operatorname{ord}_\Delta(\gamma),
 \qquad
 \operatorname{ord}_\Delta(s_\beta(\gamma))
 \le\operatorname{ord}_\Delta(\gamma),
\]
then, for every \(w\in W_{\alpha,\beta}\),
\[
 \operatorname{ord}_\Delta(w(\gamma))
 \le\operatorname{ord}_\Delta(\gamma).
\]
Moreover,
\[
 \operatorname{ord}_\Delta(w(\gamma))
 <\operatorname{ord}_\Delta(\gamma)
\]
if \(w\ne1,s_\beta\).
\end{corollary}

\begin{proof}
Consider the dual root datum \(\mathcal{R}^\ast\), and then the datum
\[
 (M^\ast,M,R'^\ast,R'),
\]
where \(R'\) is the set of roots that are rational linear
combinations of \(\alpha\) and \(\beta\). Applying
Corollary~\ref{sga3:XXI:4.3} to this datum gives the stated result.
\footnote{Editors' note: The hypotheses say that
\((\gamma,\alpha^\ast)>0\) and \((\gamma,\beta^\ast)\ge0\).
Consequently
\(\gamma-w(\gamma)\in\mathbb N\alpha+\mathbb N\beta\) for every
\(w\in W_{\alpha,\beta}\), and it is nonzero if
\(w\notin\{1,s_\beta\}\).}
\end{proof}

\section{The Weyl group: generators and relations}
\label{sga3:XXI:5}
\setcounter{theorem}{0}

Let \(\mathcal{R}\) be a root datum. Since the Weyl group is the same
for this datum and for the corresponding reduced datum, we may assume
\(\mathcal{R}\) reduced when studying the Weyl group.

Let
\[
 \Delta=\{\alpha_1,\ldots,\alpha_n\}
\]
be a system of simple roots, where \(n=\rgss(\mathcal{R})\). Let
\(n_{ij}\) be the order of \(s_{\alpha_i}s_{\alpha_j}\) in \(W\).
In particular \(n_{ii}=1\), and Remarks~\ref{sga3:XXI:3.4.2} and
\ref{sga3:XXI:3.4.3} show that the subgroup \(W_{ij}\) of \(W\)
generated by \(s_{\alpha_i}\) and \(s_{\alpha_j}\) is defined by
\[
 s_{\alpha_i}^{\,2}=s_{\alpha_j}^{\,2}
 =(s_{\alpha_i}s_{\alpha_j})^{n_{ij}}=1.
\]

\begin{theorem}\label{sga3:XXI:5.1}
The group \(W\) is the group generated by the elements
\[
 s_{\alpha_i},\qquad i=1,2,\ldots,n,
\]
subject to the relations
\[
 (s_{\alpha_i}s_{\alpha_j})^{n_{ij}}=1.
\]
\end{theorem}

\begin{proof}
We have already seen that the theorem is true when \(n=2\); we shall
use this fact in the proof. Introduce the group \(\widetilde W\)
generated by elements
\[
 T_i,\qquad i=1,2,\ldots,n,
\]
subject to the relations
\[
 (T_iT_j)^{n_{ij}}=1.
\]
In particular \(n_{ii}=1\), whence \(T_i^2=1\). Let
\[
 p:\widetilde W\longrightarrow W
\]
be the group homomorphism sending \(T_i\) to \(s_{\alpha_i}\). We know
that \(p\) is surjective; we shall prove that it is injective.
\end{proof}

\begin{lemma}\label{sga3:XXI:5.2}
For every \(\alpha\in\mathscr{P}(\Delta)\), one can define uniquely an
element \(T_\alpha\in\widetilde W\) with the following properties:
\begin{enumerate}
\item[(i)] \(p(T_\alpha)=s_\alpha\);
\item[(ii)] \(T_{\alpha_i}=T_i\);
\item[(iii)] if \(\beta\) and \(\alpha\) are positive roots such that
\(s_{\alpha_i}(\alpha)=\beta\), then
\[
 T_iT_\alpha T_i=T_\beta.
\]
\end{enumerate}
\end{lemma}
